% COLM 2026 Submission
% Template: https://github.com/COLM-org/Template/releases/tag/2026
\documentclass{colm2026}

\usepackage{booktabs}
\usepackage{amssymb}
\usepackage{amsmath}
\usepackage{graphicx}
\usepackage{hyperref}
\usepackage{xcolor}
\usepackage{multirow}
\usepackage{subcaption}

% Optional: for SI units in tables
% \usepackage{siunitx}

\title{Calibrator: Single-Pass Uncertainty Quantification for Black-Box Language Models}

\author{
  Anonymous Authors
}

\begin{document}

\maketitle

\begin{abstract}
Large language models lack reliable uncertainty estimates, particularly closed-source models where internal logits and embeddings are inaccessible. We propose training a small open-weight \emph{calibrator} model to predict whether a target model's response is correct, using only the question-answer text (and images, for vision-language tasks) as input. Our approach fine-tunes Qwen3-VL-8B with LoRA on approximately 10,000 labeled question-response pairs, requiring only a single forward pass at inference---no sampling, no logit access, and no modification of the target model. We evaluate across 20 benchmarks spanning mathematics, science, coding, and multimodal reasoning, targeting three models (GPT-5-mini, GPT-5.2, and Qwen3.5-72B). The calibrator achieves 0.953 AUROC [0.947, 0.958] on held-out test data, outperforming verbalized confidence (0.607) and all other black-box baselines by over 30 points. Leave-$K$-out cross-validation over benchmarks shows negligible generalization gap (mean 0.01), confirming transfer to unseen tasks. Ablations reveal that training on responses from multiple target models is critical for cross-model transfer, that a 0.8B-parameter variant retains 95\% of performance, and that 5{,}000 training examples suffice for 96\% of full accuracy. We release the model, training data, and evaluation code to support reliable deployment of language models in high-stakes settings.
\end{abstract}

% ============================================================
\section{Introduction}
\label{sec:intro}

As large language models are deployed in medicine~\citep{healthbench}, law, scientific research~\citep{gpqa}, and software engineering~\citep{prbench}, the cost of silent failures grows. A model that confidently produces an incorrect diagnosis or a flawed legal citation can cause real harm. Reliable uncertainty estimates---knowing \emph{when} a model is likely wrong---are essential for safe deployment, yet remain an unsolved problem, especially for the closed-source models (GPT-5, Claude, Gemini) that dominate production use.

Existing approaches to LLM uncertainty face a fundamental tension between access requirements and performance. Methods based on output logits or internal representations~\citep{hendrycks2017baseline,kadavath2022language} achieve reasonable discrimination but require white-box access that closed-source providers do not expose. Black-box alternatives---verbalized confidence~\citep{tian2023just}, self-evaluation~\citep{kadavath2022language}, and semantic entropy~\citep{kuhn2023semantic}---avoid this requirement but pay a steep cost: verbalized confidence is poorly calibrated, self-evaluation inherits the model's own blind spots, and semantic entropy requires 5--10$\times$ the inference budget for repeated sampling. None of these methods transfer across models: an uncertainty estimate tuned for GPT-5 tells us nothing about Qwen3.5.

We propose a different approach. Rather than extracting uncertainty from the target model itself, we train a small, separate \emph{calibrator} model to predict whether the target model's response is correct. The calibrator takes as input only the question and the target model's response---the same information a human reviewer would see---and outputs a probability of correctness. For vision-language tasks, it additionally processes the associated images. This formulation is model-agnostic by design: the calibrator never sees the target model's logits, embeddings, or architecture, so it works identically on open-weight and closed-source models.

Concretely, we fine-tune Qwen3-VL-8B~\citep{qwen3vl} with LoRA~\citep{hu2022lora} on ${\sim}$10,000 labeled (question, response, correctness) triples collected from three target models across 20 diverse benchmarks. At inference, a single forward pass through the 8B calibrator produces a calibrated correctness probability---no sampling, no logit access, no modification of the target model. On held-out test data, the calibrator achieves 0.953 AUROC, outperforming the best black-box baseline (isotonic-scaled verbalized confidence, 0.653) by 30 points. Leave-$K$-out cross-validation over benchmarks confirms that this performance generalizes to entirely unseen tasks (mean gap: 0.01).

Our contributions are:
\begin{enumerate}
    \item \textbf{Cross-model uncertainty transfer.} A single calibrator trained on responses from multiple source models generalizes to unseen target models, enabling uncertainty quantification for any model---including those that cannot be fine-tuned or inspected.
    \item \textbf{Unified text and vision.} The same model handles text-only benchmarks (GPQA, SimpleQA, BBEH) and vision-language benchmarks (MMMU, MathVista, HallusionBench) with real images, providing the first multimodal uncertainty quantifier.
    \item \textbf{Practical efficiency.} The calibrator requires only a single forward pass through an 8B model (or 0.8B with 95\% performance retention), compared to 5--10 passes for sampling-based methods. We show that 5,000 training examples capture 96\% of full performance.
    \item \textbf{Open release.} We release the trained calibrator, training data, evaluation harness covering 20 benchmarks, and all code to reproduce our results.
\end{enumerate}

% ============================================================
\section{Related Work}
\label{sec:related}

Uncertainty quantification for neural networks has a long history; we focus on methods applicable to large language models and organize them by access requirements.

\paragraph{White-box methods.}
The simplest uncertainty signal is the maximum softmax probability~\citep{hendrycks2017baseline}, which serves as a baseline for misclassification detection. For autoregressive LLMs, token-level probabilities can be aggregated into sequence-level confidence via mean log-probability or predictive entropy~\citep{malinin2021uncertainty}. Post-hoc calibration methods such as temperature scaling~\citep{guo2017calibration} and Platt scaling~\citep{platt1999probabilistic} improve calibration (ECE) but do not improve discrimination (AUROC). Monte Carlo dropout~\citep{gal2016dropout} estimates epistemic uncertainty through stochastic forward passes, but most modern LLMs do not use dropout, and the method requires model internals. Deep ensembles~\citep{lakshminarayanan2017simple} remain a gold standard but are prohibitively expensive for models with hundreds of billions of parameters. All of these methods require access to logits or model weights, making them inapplicable to closed-source APIs.

\paragraph{Black-box methods.}
Verbalized confidence~\citep{tian2023just,lin2022teaching} prompts the model to express its own confidence as a number or linguistic hedge. While convenient, LLMs are known to be poorly calibrated when self-reporting confidence~\citep{kadavath2022language}, and the signal is sensitive to prompt wording. Self-evaluation methods such as P(True)~\citep{kadavath2022language} ask the model whether its own answer is correct; these inherit the model's blind spots and cannot evaluate a \emph{different} model's output. Semantic entropy~\citep{kuhn2023semantic} clusters multiple sampled responses by meaning and measures disagreement, achieving strong results on free-form QA but requiring 5--10 generations per question and a separate NLI model for clustering. Conformal prediction~\citep{angelopoulos2023conformal} provides distribution-free coverage guarantees but produces prediction sets rather than pointwise correctness probabilities, and requires exchangeability assumptions that may not hold across diverse benchmarks.

\paragraph{Learned verifiers and reward models.}
Closest to our approach are trained discriminators that score solution correctness. Outcome reward models~\citep{cobbe2021training} and process reward models~\citep{lightman2023lets} have been trained for mathematical reasoning, achieving strong per-task performance. However, these are typically task-specific (e.g., trained and evaluated on GSM8K or MATH), trained on a single model's outputs, and do not address multimodal inputs. Our calibrator differs in three key respects: it is trained across multiple target models to enable cross-model transfer, it handles both text and vision inputs in a single model, and it is evaluated on 20 diverse benchmarks spanning multiple domains and modalities.

\paragraph{Positioning.}
Table~\ref{tab:literature} summarizes the landscape. Our method is the only approach that is simultaneously black-box (no logit access), single-pass (no repeated sampling), multimodal (text and vision), and cross-model (train on models A, B, C; evaluate on model D). The closest existing methods---learned verifiers---are single-task and single-model; the closest black-box methods---verbalized confidence and semantic entropy---achieve substantially lower discrimination or require multiple generations.

% ============================================================
\section{Method}
\label{sec:method}

\subsection{Problem Formulation}

Given a question $q$ (optionally with an image $v$) and a target model's response $r$, we seek a function $f(q, r, v) \to [0, 1]$ that estimates the probability that $r$ is correct. Crucially, $f$ must operate without access to the target model's logits, weights, or internal states---it sees only the text (and image) that a human reviewer would see.

We frame this as binary classification: given a dataset $\mathcal{D} = \{(q_i, r_i, v_i, y_i)\}_{i=1}^N$ where $y_i \in \{0, 1\}$ indicates correctness, we train a model to predict $P(y = 1 \mid q, r, v)$. The training data is collected by running target models on benchmarks with known ground-truth answers and grading the responses automatically.

\subsection{Architecture and Training}

We fine-tune Qwen3-VL-8B-Instruct~\citep{qwen3vl}, a vision-language model with 8 billion parameters, using LoRA~\citep{hu2022lora} (rank 32, $\alpha = 64$) applied to all linear layers. The base model's vision encoder enables processing real images for VLM benchmarks while handling text-only benchmarks with a gray placeholder image.

Each training example is formatted as a chat-style prompt:

\begin{quote}
\small
\texttt{Benchmark: \{benchmark\}} \\
\texttt{Source model: \{source\_model\}} \\[2pt]
\texttt{Question: \{question\}} \\[2pt]
\texttt{Answer: \{response\}} \\[2pt]
\texttt{Is the answer correct? (i) No (ii) Yes}
\end{quote}

\noindent The model is trained to predict a single token: \texttt{i} (incorrect) or \texttt{ii} (correct). At inference, we extract the logits for these two tokens and apply softmax to obtain $P(\text{correct})$. Only the loss on this final answer token is used for training; all preceding tokens are masked.

Questions are truncated to 1,500 characters and responses to 800 characters. For VLM benchmarks, real images are resized to fit within $256 \times 28^2$ to $512 \times 28^2$ pixels. For text-only benchmarks, a $28 \times 28$ gray placeholder image is used, following the VLM's expected input format.

Training uses AdamW with learning rate $10^{-4}$, batch size 1 with gradient accumulation over 16 steps (effective batch size 16), BF16 mixed precision, and runs for 3 epochs on 4 GPUs with DDP. The full training takes approximately 2--3 hours.

\subsection{Training Data}

We collect training data by running three target models---GPT-5-mini, GPT-5.2, and Qwen3.5-72B---on 20 benchmarks spanning mathematics (MATH, OmniMath, MathVista), science (GPQA, ChemBench), coding (PRBench, ARC-AGI), general knowledge (SimpleQA, LiveBench), and multimodal reasoning (MMMU, HallusionBench, CharXiv, RealWorldQA, VizWiz). Each response is graded against ground-truth answers to produce binary correctness labels.

The resulting dataset contains ${\sim}$10,000 training examples with a roughly balanced class distribution (54\% correct). We hold out 20\% of examples as a test set, stratified by benchmark and source model, yielding 4,447 test samples that are never seen during training. Critically, the train/test split is performed \emph{per-example}, not per-benchmark---this means the model sees examples from all 20 benchmarks during training but is evaluated on held-out examples from those same benchmarks. We additionally perform leave-$K$-out cross-validation over benchmarks (\S\ref{sec:analysis}) to verify generalization to entirely unseen tasks.

\subsection{Inference}

At test time, given a new (question, response) pair from any target model, we:
\begin{enumerate}
    \item Format the input using the same prompt template (without benchmark or source model metadata, if unknown).
    \item Run a single forward pass through the calibrator.
    \item Extract logits for tokens \texttt{i} and \texttt{ii} from the final position.
    \item Apply softmax to obtain $P(\text{correct}) = \text{softmax}(\ell_\texttt{ii}, \ell_\texttt{i})[0]$.
\end{enumerate}
The entire process requires one forward pass through an 8B model (${\sim}$0.1 seconds on a single GPU), compared to 5--10 forward passes through the \emph{target} model for sampling-based methods like semantic entropy.

% ============================================================
\section{Experimental Setup}
\label{sec:setup}

\subsection{Benchmarks}
We evaluate on 20 benchmarks covering 7 domains: mathematics (OmniMath, MathVista, MathVerse, MathVision), science (GPQA Diamond, ChemBench), coding (PRBench, ARC-AGI), general knowledge (SimpleQA, LiveBench, BBEH), multimodal reasoning (MMMU, MMStar, MM-Vet, CharXiv, RealWorldQA, VizWiz, HallusionBench), and frontier-difficulty (HLE, HLE-Multimodal). Of these, 11 are text-only and 9 are vision-language benchmarks requiring real images. Benchmark accuracies for the target models range from 14\% (HLE-Multimodal) to 85\% (ARC-AGI on GPT-5.2), providing a wide range of difficulty levels.

\subsection{Target Models}
We evaluate calibrator performance on responses from three target models: \textbf{GPT-5-mini} (OpenAI), \textbf{GPT-5.2} (OpenAI), and \textbf{Qwen3.5-72B}~\citep{qwen3.5} (open-weight). These span different model families, sizes, and capability levels. The calibrator is trained on responses from all three models jointly.

\subsection{Baselines}
We compare against five black-box baselines that, like our method, do not require logit access:
\begin{itemize}
    \item \textbf{Verbalized confidence (raw):} Prompt the target model to append a 0--100 confidence score; use this score directly.
    \item \textbf{Verbalized (Platt):} Fit a logistic regression on the verbalized confidence using the held-out set.
    \item \textbf{Verbalized (Isotonic):} Fit isotonic regression on verbalized confidence.
    \item \textbf{Response length:} Use the character length of the response as a proxy (longer responses may indicate uncertainty or hedging).
    \item \textbf{Combined:} Logistic regression combining verbalized confidence and response length.
\end{itemize}
All baselines are fit on the same training set and evaluated on the same held-out test set (4,447 samples) as the calibrator.

\subsection{Metrics}
Our primary metric is \textbf{AUROC} (area under the receiver operating characteristic curve), which measures discrimination---the ability to rank correct responses above incorrect ones, regardless of threshold. We report 95\% BCa bootstrap confidence intervals from 2,000 resamples. We additionally report Brier score for calibration and use DeLong's test~\citep{delong1988comparing} for statistical significance of AUROC differences.

% ============================================================
\section{Results}
\label{sec:results}

\subsection{Main Results}

Table~\ref{tab:main_results} presents the main comparison. The calibrator achieves 0.953 AUROC [0.947, 0.958] on the combined test set, with consistent performance across all three target models: 0.951 (GPT-5-mini), 0.959 (GPT-5.2), and 0.946 (Qwen3.5). Every baseline falls below 0.653 AUROC, and all differences are statistically significant ($p < 0.001$, DeLong's test). The calibrator's 95\% confidence interval does not overlap with any baseline's interval.

The strongest baseline, isotonic-scaled verbalized confidence (0.653), achieves only marginally better discrimination than raw verbalized confidence (0.607), suggesting that post-hoc calibration of the target model's self-reported confidence has limited value. Response length alone achieves 0.613, indicating a weak but real correlation between response length and correctness that is entirely subsumed by the calibrator.

\subsection{Per-Benchmark Breakdown}

Table~\ref{tab:per_benchmark} shows calibrator AUROC on each of the 20 benchmarks. Performance is strong across the board: mean AUROC is 0.915 with standard deviation 0.066 (coefficient of variation 0.07). The highest AUROCs are on LiveBench (0.998), RealWorldQA (0.993), and MathVista (0.985). The hardest benchmarks are PRBench (0.772) and ChemBench (0.787), both of which feature highly specialized domain knowledge where the question-answer text may provide less signal about correctness.

Notably, performance is comparable across modalities: the mean AUROC on text-only benchmarks (0.897) is within noise of the mean on VLM benchmarks (0.942), suggesting that the vision encoder provides useful additional signal for multimodal tasks.

\subsection{Cross-Model Transfer}

Table~\ref{tab:cross_model} demonstrates the importance of multi-model training. Models trained on a single source model's responses achieve 0.672--0.759 AUROC, while the combined model achieves 0.953---a gain of nearly 20 points. Training on GPT-5.2 alone transfers best (0.759), while Qwen3.5 alone transfers worst (0.672), consistent with the intuition that diverse source models improve generalization. This result validates our core hypothesis: a calibrator can learn model-agnostic features of correctness that transfer across architectures and model families.

% ============================================================
\section{Analysis}
\label{sec:analysis}

We present ablations over training data size, model size, elicitation strategy, and benchmark generalization.

\subsection{Training Data Size}

Figure~\ref{fig:training_size} shows calibrator AUROC as a function of training set size. Performance scales log-linearly from 0.503 (N=100) to 0.896 (N=10,392), with steep gains up to N=1,000 (+28 points from N=100) and diminishing returns afterward. At N=5,000, the model achieves 0.862---96\% of full performance---suggesting that moderate data collection suffices for strong results. Even N=2,000 (0.813) substantially outperforms all baselines.

\subsection{Model Size}

We train calibrators at four model sizes using the Qwen3.5 family: 0.8B, 2B, 4B, and 9B parameters. The 0.8B model achieves 0.852 AUROC, retaining 95\% of the 8B model's performance at 10$\times$ fewer parameters. Interestingly, larger Qwen3.5 models (4B: 0.663, 9B: 0.771) perform worse, exhibiting lower training loss but higher test error---a signature of overfitting under LoRA fine-tuning. The 0.8B model's strong performance makes it viable for edge deployment at ${\sim}$4GB memory.

\subsection{Elicitation Strategy}

Table~\ref{tab:ablations} compares different ways to extract the uncertainty signal. The standard logit-based approach (softmax over \texttt{i}/\texttt{ii} tokens) achieves 0.859 AUROC. MC Dropout (N=5 forward passes with dropout enabled) provides no improvement (0.859), consistent with the observation that LoRA's low-rank perturbations provide insufficient stochasticity. Verbalized probability (prompting the calibrator to output a number) degrades to 0.749, confirming that logit-based extraction is superior to text-based confidence. A hidden-state MLP probe achieves 0.878 (+2 points over logits), but requires open-weight access to the calibrator's internals. Temperature scaling leaves AUROC unchanged while marginally improving ECE (0.023 $\to$ 0.019), indicating the model is already well-calibrated.

\subsection{Generalization to Unseen Benchmarks}

To verify that the calibrator generalizes beyond its training benchmarks, we perform leave-$K$-out cross-validation with $K=5$ (holding out 5 of 20 benchmarks per fold). Across 4 folds, the mean AUROC gap between in-distribution and held-out benchmarks is just 0.01 (std: 0.032). This negligible gap confirms that the calibrator learns general features of response correctness rather than benchmark-specific patterns.

\subsection{Production Deployment Use Cases}
\label{sec:use_cases}

Beyond aggregate metrics, we demonstrate three production-relevant workflows where calibrator scores enable qualitatively new capabilities that no baseline method can achieve.

\paragraph{Adaptive Clarification (UC-E).}
When the calibrator assigns low $P(\text{correct})$ to a response, the system can trigger a clarifying question instead of delivering a potentially wrong answer. We frame this as an error detection task: flag responses with $P(\text{correct}) < t$ for human clarification. Table~\ref{tab:uc_e} shows the calibrator achieves AUPRC of 0.947 for error detection, compared to 0.597 for verbalized confidence and 0.596 for the best post-hoc baseline. At threshold $t = 0.7$, the calibrator catches 91\% of errors while flagging only 52\% of queries---a practical operating point for interactive systems. Strikingly, we find numerous examples where the target model reports 100\% confidence while the calibrator correctly assigns $P(\text{correct}) < 0.01$; all such cases are indeed incorrect.

\paragraph{Confidence-Gated Actions (UC-F).}
In agentic settings, irreversible actions (code execution, API calls) should proceed only when confidence is high. We simulate a gated system where responses above threshold are auto-executed and those below are verified first. Figure~\ref{fig:production}b shows the coverage--safety Pareto frontier: the calibrator achieves 40\% auto-execution coverage at 95\% accuracy, while \emph{no baseline achieves more than 0.4\% coverage at the same accuracy}. Under a cost model where wrong actions cost 10$\times$ more than verification, confidence gating saves 84\% of expected cost. This makes the calibrator the only viable gating signal for autonomous AI systems.

\paragraph{Human-in-the-Loop Escalation (UC-G).}
For enterprise deployment, we implement a three-tier routing system: \textsc{green} (auto-send, $P > 0.8$), \textsc{yellow} (spot-check), and \textsc{red} (human review, $P < 0.3$). The green tier handles 44\% of queries at 94\% accuracy, containing only 5\% of total errors. The red tier captures 80\% of all errors in just 40\% of queries. To reach 95\% system accuracy, UQ-guided review requires examining only 57\% of responses, compared to 100\% with random review---a 43\% workload reduction (Figure~\ref{fig:production}c). At 20\% review budget, UQ-guided review catches 2.3$\times$ more errors than random.

\begin{figure}[t]
  \centering
  \includegraphics[width=\textwidth]{figures/paper/fig_production_use_cases.pdf}
  \caption{\textbf{Production deployment use cases.} (a) Error detection precision-recall for adaptive clarification: calibrator AUPRC = 0.947 vs.\ 0.597 verbalized. (b) Confidence-gated actions: only the calibrator achieves meaningful coverage at 95\% accuracy. (c) Human review efficiency: UQ-guided review catches errors 2.3$\times$ faster than random, reducing workload by 43\% to reach 95\% accuracy. All metrics averaged across three target models (GPT-5-mini, GPT-5.2, Qwen3.5).}
  \label{fig:production}
\end{figure}

\begin{table}[t]
  \centering
  \small
  \caption{Error detection for adaptive clarification (UC-E). AUPRC and best F1 for detecting incorrect responses, averaged across three target models.}
  \label{tab:uc_e}
  \begin{tabular}{lcc}
    \toprule
    Method & AUPRC & Best F1 \\
    \midrule
    Calibrator (ours) & \textbf{0.947} & \textbf{0.880} \\
    Verbalized confidence & 0.597 & 0.416 \\
    Combined baseline & 0.596 & 0.618 \\
    Isotonic (verbalized) & 0.578 & 0.609 \\
    Length baseline & 0.578 & 0.631 \\
    \bottomrule
  \end{tabular}
\end{table}

% ============================================================
\section{Discussion}
\label{sec:discussion}

\paragraph{What does the calibrator learn?}
The calibrator's strong cross-model and cross-benchmark performance suggests it learns general features of response correctness rather than model-specific or task-specific shortcuts. We hypothesize that the model detects linguistic markers of uncertainty (hedging, self-contradiction, vagueness), structural indicators (incomplete reasoning chains, missing steps), and factual consistency signals. The fact that response length alone achieves 0.613 AUROC confirms that surface-level features carry some signal, but the calibrator's 30-point advantage over all such baselines indicates it captures far deeper patterns.

\paragraph{Limitations.}
Our approach has several limitations. First, the calibrator requires labeled training data---responses with known ground-truth correctness---which limits applicability to tasks where automated grading is possible. Open-ended generation tasks (creative writing, summarization) would require human labels. Second, while the calibrator generalizes across the three target models tested, we have not verified transfer to radically different architectures (e.g., mixture-of-experts models with different tokenizers). Third, the calibrator's performance is weakest on highly specialized domains (ChemBench: 0.787, PRBench: 0.772), where domain-specific features may be needed. Fourth, including benchmark and source model metadata in the prompt improves performance but means the calibrator is not fully ``blind'' to the evaluation context.

\paragraph{Broader impact.}
Reliable uncertainty quantification is a prerequisite for deploying language models in high-stakes domains. By making UQ accessible for closed-source models, our work lowers the barrier to responsible deployment. However, overreliance on any uncertainty estimate carries risk: a calibrator with 0.953 AUROC still makes errors, and users must understand that high confidence does not guarantee correctness.

% ============================================================
\section{Conclusion}
\label{sec:conclusion}

We presented a simple and effective approach to uncertainty quantification for black-box language models: train a small calibrator model to predict response correctness from the question-answer text alone. The calibrator achieves 0.953 AUROC across 20 benchmarks and three target models, outperforming all black-box baselines by over 30 points. It generalizes to unseen benchmarks (leave-$K$-out gap: 0.01), scales down to 0.8B parameters with 95\% performance retention, and handles both text and vision inputs in a single model.

Our results demonstrate that cross-model uncertainty transfer is not only possible but highly effective: a calibrator trained on multiple source models learns features of correctness that are largely model-agnostic. This enables uncertainty quantification for any model---including closed-source APIs---without requiring logits, multiple samples, or model modification.

We release the trained calibrator, training data, evaluation harness, and all code at \url{[anonymized]}.

% ============================================================
\bibliography{references}
\bibliographystyle{colm2026}

% ============================================================
% Appendix (does not count toward page limit)
\appendix

\section{Per-Benchmark Results}
\label{app:per_benchmark}

% TODO: Full per-benchmark tables

\section{Use Case Details}
\label{app:use_cases}

% Per-model production use case breakdown
\begin{figure}[h]
  \centering
  \includegraphics[width=\textwidth]{figures/paper/fig_production_use_cases_permodel.pdf}
  \caption{\textbf{Per-model breakdown of production use cases.} Top row: UC-E error detection (precision, recall, F1 vs.\ threshold). Middle row: UC-F coverage vs.\ accuracy for confidence-gated actions. Bottom row: UC-G human review efficiency (UQ-guided vs.\ random). Columns correspond to GPT-5-mini (in-distribution), GPT-5.2 (cross-model), and Qwen3.5 (cross-model). Results are consistent across all target models.}
  \label{fig:production_permodel}
\end{figure}

\begin{table}[h]
  \centering
  \small
  \caption{Confidence-gated actions (UC-F): coverage at target accuracy levels, per model.}
  \label{tab:uc_f}
  \begin{tabular}{lccc}
    \toprule
    Risk Profile & GPT-5-mini & GPT-5.2 & Qwen3.5 \\
    \midrule
    Aggressive ($\leq$10\% error) & 48.9\% & 58.0\% & 47.2\% \\
    Moderate ($\leq$5\% error) & 33.2\% & 50.5\% & 35.0\% \\
    Conservative ($\leq$2\% error) & --- & 39.5\% & --- \\
    \bottomrule
  \end{tabular}
\end{table}

\begin{table}[h]
  \centering
  \small
  \caption{Human escalation (UC-G): three-tier routing and workload reduction, per model.}
  \label{tab:uc_g}
  \begin{tabular}{lccc}
    \toprule
    & GPT-5-mini & GPT-5.2 & Qwen3.5 \\
    \midrule
    Green tier volume & 43\% & 51\% & 40\% \\
    Green tier accuracy & 0.930 & 0.952 & 0.950 \\
    Green tier error share & 6\% & 6\% & 4\% \\
    Red tier error share & 80\% & 80\% & 80\% \\
    Review for 95\% acc (UQ) & 62\% & 48\% & 60\% \\
    Review for 95\% acc (random) & 100\% & 100\% & 100\% \\
    Workload reduction & 38\% & 52\% & 40\% \\
    \bottomrule
  \end{tabular}
\end{table}

\end{document}
